\subsection{多项式环上分次模的庞加莱级数}

令 \(H_{0}\) 为一个环，I 为有限集，H 为多项式环 \(\mathrm{H}_{0}\big[(\mathrm{X}_{i})_{i\in\mathbf{I}}\big]\)。对于每个 \(i\in I\)，令 \(d_{i}\) 为整数 \(>0\)。赋予 H 一个 Z 型分次环结构，使 \(H_{0}\) 的元素都是 0 次齐次的，而每个 \(X_{i}\) 都是 \(d_{i}\) 次齐次的。当 \(d_{i}=1\) 对每个 \(i\) 都成立时，就得到多项式环的通常分次。

令 M 为有限生成的分次 H-模，且其所有齐次分量都是有限长度的 \(H_{0}\)-模；称 M 的庞加莱级数为元素 \(P_{M}\)，它属于 \(\mathbf{Z}((T))\) 且满足 \(\mathrm{P}_{\mathrm{M}} = \sum_{n \in \mathbf{Z}} \mathrm{long}_{\mathrm{H}_{0}}(\mathrm{M}_{n}) \cdot \mathrm{T}^{n}\)，并置 \(Q_{M} = P_{M} \cdot \prod_{i \in I} (1 - T^{d_{i}})\)。

定理 1. --- \(\mathbf{Q}_{\mathbf{M}}\) 是 \(\mathbf{Z}((\mathbf{T}))\) 的元素，且属于 \(\mathbf{Z}[\mathbf{T}, \mathbf{T}^{-1}]\)。

将 \(H_{0}\) 除以 M 作为 \(H_{0}\)-模的湮灭子，我们归约到 M 是忠实 \(H_{0}\)-模的情形。若 \(a, b \in Z\) 使得 M 作为 H-模由 \(M' = \sum_{a \leqslant i \leqslant b} M_{i}\) 生成，则 \(M'\) 是非零、有限生成、忠实且具有有限长度的 \(H_{0}\)-模；因此，

环 \(\mathbf{H}_0\) 是阿廷环（A, VIII, § 1, no. 3），因而是诺特环（loc. cit., § 9, no. 1）。所以多项式环 H 也是诺特环（loc. cit., § 1, no. 4）。若 I 为空，则 \(\mathbf{H} = \mathbf{H}_0\)，且整数族 \((\mathrm{long}_{\mathbf{H}_0}(\mathbf{M}_n))_{n \in \mathbf{Z}}\) 的支集有限，因为 M 是有限生成的 \(\mathbf{H}_0\)-模；故定理在此情形成立。现在对集合 I 的基数归纳论证，设 I 非空；取 \(j \in \mathbf{I}\)，并令 \(\mathbf{J} = \mathbf{I} - \{j\}\)。记 H' 为由 \(\mathbf{H}_0\) 以及 \(X_i\)（其中 \(i\) 属于 J）生成的 H 的分次子环；考虑 M 上的倍乘映射 \((\mathbf{X}_j)_\mathbf{M}\)，其比率为 \(X_j\)，以及它的核 R 和余核 S。对于每个 \(n \in \mathbf{Z}\)，有 \(\mathbf{H}_0\)-模的正合序列

\[
0 \rightarrow \mathrm{R} _ {n - d _ {j}} \rightarrow \mathrm{M} _ {n - d _ {j}} \rightarrow \mathrm{M} _ {n} \rightarrow \mathrm{S} _ {n} \rightarrow 0,
\]

因而 \(\mathbf{R}_{n - d_j}\) 和 \(\mathbf{S}_n\) 具有有限长度，并且有

\[
\operatorname{long} _ {\mathrm{H} _ {0}} \left(\mathrm{M} _ {n}\right) - \operatorname{long} _ {\mathrm{H} _ {0}} \left(\mathrm{M} _ {n - d _ {j}}\right) = \operatorname{long} _ {\mathrm{H} _ {0}} \left(\mathrm{S} _ {n}\right) - \operatorname{long} _ {\mathrm{H} _ {0}} \left(\mathrm{R} _ {n - d _ {j}}\right).\tag{4}
\]

由于 M 是诺特环 H 上的有限生成模，H-模 R 和 S 都是有限生成的；由于它们被 \(X_{j}\) 湮灭，所以它们是有限生成的 H'-模。根据归纳假设，\(P_{R}\cdot\prod_{i\in J}(1-T^{d_{i}})\) 和 \(P_{S}\cdot\prod_{i\in J}(1-T^{d_{i}})\) 这两个 \(\mathbf{Z}((T))\) 的元素属于 \(\mathbf{Z}[T,T^{-1}]\)；由 (4) 得到

\[
\mathrm{P} _ {\mathrm{M}} - \mathrm{T} ^ {d _ {j}}. \mathrm{P} _ {\mathrm{M}} = \mathrm{P} _ {\mathrm{S}} - \mathrm{T} ^ {d _ {j}}. \mathrm{P} _ {\mathrm{R}},
\]

即 \((1 - \mathrm{T}^{d_j}).\mathrm{P}_{\mathrm{M}} = \mathrm{P}_{\mathrm{S}} - \mathrm{T}^{d_j}.\mathrm{P}_{\mathrm{R}}\)；于是有

\[
\mathrm{P} _ {\mathrm{M}} \cdot \prod_ {i \in \mathrm{I}} (1 - \mathrm{T} ^ {d _ {i}}) = \mathrm{P} _ {\mathrm{S}} \cdot \prod_ {i \in \mathrm{J}} (1 - \mathrm{T} ^ {d _ {i}}) - \mathrm{T} ^ {d _ {j}} \cdot \mathrm{P} _ {\mathrm{R}} \cdot \prod_ {i \in \mathrm{J}} (1 - \mathrm{T} ^ {d _ {i}}),\tag{5}
\]

从而得到结论。

例 1. --- 假设 \(H_{0}\) 是阿廷环，并取 M = H。用前面的记号，有 R = 0 且 S = H'，所以由 (5) 有 \(Q_{H} = Q_{H'}\)；由于 \(Q_{H_{0}} = \text{long}(H_{0})\)，由归纳可得 \(Q_{H} = \text{long}(H_{0})\)，即

\[
\mathrm{P} _ {\mathrm{H}} = \operatorname{long} (\mathrm{H} _ {0}). \prod_ {i \in \mathrm{I}} (1 - \mathrm{T} ^ {d _ {i}}) ^ {- 1}.\tag{6}
\]

以下设 H 取通常的分次，其中 \(d_{i}=1\) 对每个 \(i\in I\) 都成立，并置 \(r=\operatorname{Card}(\mathbf{I})\)；有 \(\mathrm{P}_{\mathrm{M}}=\mathrm{Q}_{\mathrm{M}}(\mathrm{T}).(1-\mathrm{T})^{-r}\)。置 \(c_{\mathrm{M}}=\mathrm{Q}_{\mathrm{M}}(1)\)。于是由 no. 1 的公式 (2)，有：

推论。--- a) 若 r = 0，则有 long \(_{H_{0}}(M)\) = c \(_{M}\)。

\begin{enumerate}
\def\labelenumi{\alph{enumi})}
\setcounter{enumi}{1}
\item
  若 \(r = 1\)，则有 \(\mathrm{long}_{\mathrm{H}_0}(\mathbf{M}_n) = c_{\mathbf{M}}\)，当 \(n\) 充分大时。
\item
  若 \(r > 1\)，则有 long \(_{\mathrm{H}_0}(\mathbf{M}_n) = c_{\mathbf{M}} \frac{n^{r-1}}{(r - 1)!} + \rho_n n^{r-2}\)，其中 \(\rho_n\) 趋于 \(\mathbf{R}\) 中的一个极限，当 \(n\) 无限增大时。
\end{enumerate}

备注。--- 1) 根据引理 2，整数 \(c_{\mathbf{M}}\) 为正。可能有 \(\mathbf{M} \neq 0\) 而 \(c_{\mathbf{M}} = 0\)（cf. prop. 2）。

\begin{enumerate}
\def\labelenumi{\arabic{enumi})}
\setcounter{enumi}{1}
\tightlist
\item
  设 \(0 \to M' \to M \to M'' \to 0\) 是分次 H-模及 0 次同态的正合序列，且 M 是 H 上有限生成的，并且 \(M_n\) 都是 \(H_0\) 上具有有限长度的模，对于每个 \(n\)。则对于每个 \(n \in \mathbf{Z}\)，有
\end{enumerate}

\[
\operatorname{long} _ {\mathrm{H} _ {0}} \left(\mathrm{M} _ {n}\right) = \operatorname{long} _ {\mathrm{H} _ {0}} \left(\mathrm{M} _ {n} ^ {\prime}\right) + \operatorname{long} _ {\mathrm{H} _ {0}} \left(\mathrm{M} _ {n} ^ {\prime \prime}\right),
\]

因此 \(\mathrm{P_M} = \mathrm{P_{M'}} + \mathrm{P_{M''}}\)、\(\mathrm{Q_M} = \mathrm{Q_{M'}} + \mathrm{Q_{M''}}\)，且 \(c_{\mathbf{M}} = c_{\mathbf{M'}} + c_{\mathbf{M''}}\)。

\begin{enumerate}
\def\labelenumi{\arabic{enumi})}
\setcounter{enumi}{2}
\tightlist
\item
  设 \(\mathbf{M}(p)\) 是由 \(\mathbf{M}\) 在分次中平移 \(p\) 得到的模（A, II, p. 165, example 3）。由于有 \(\mathbf{M}(p)_n = \mathbf{M}_{p+n}\)，所以 \(\mathbf{P}_{\mathbf{M}(p)} = \mathrm{T}^{-p}\mathbf{P}_{\mathbf{M}}\)、\(\mathbf{Q}_{\mathbf{M}(p)} = \mathrm{T}^{-p}\mathbf{Q}_{\mathbf{M}}\)，且 \(c_{\mathbf{M}(p)} = c_{\mathbf{M}}\)。
\end{enumerate}

例。--- 2) 假设 H \(_{0}\) 是阿廷环，且 M 是由 s 个齐次、线性无关的元素生成的分次自由 H-模，这些元素的次数分别为 \(\delta_{1}, ..., \delta_{s}\)。则 M 同构于 H( \(-\delta_{1}\) ) ⊕ \(\cdots\) ⊕ H( \(-\delta_{s}\) )。因此由备注 2、3 和例 1，有

\[
\begin{array}{l} \mathrm {P_ {M}} = \operatorname{long} (\mathrm {H_ {0}}) \left(\sum_ {i = 1} ^ {s} \mathrm{T} ^ {\delta_ {i}}\right) (1 - \mathrm{T}) ^ {- r}, \\ \mathrm {Q_ {M}} = \operatorname{long} (\mathrm {H_ {0}}) \left(\sum_ {i = 1} ^ {s} \mathrm{T} ^ {\delta_ {i}}\right), \\ c _ {\mathrm{M}} = s. \operatorname{long} (\mathrm {H_ {0}}). \end{array}
\]

\begin{enumerate}
\def\labelenumi{\arabic{enumi})}
\setcounter{enumi}{2}
\tightlist
\item
  再次假设 \(H_{0}\) 是阿廷环；令 M 为分次 H-模，并假设存在分次 H-模及 0 次同态的正合序列
\end{enumerate}

\[
0 \to \mathrm{L} _ {n} \to \mathrm{L} _ {n - 1} \to \dots \to \mathrm{L} _ {0} \to \mathrm{M} \to 0,
\]

其中，对于 \(k=0,1,...,n\)，\(\mathbf{L}_{k}\) 是由线性无关的齐次元素生成的分次自由 H-模，这些元素的次数分别为 \(\delta_{k,1},..., \delta_{k,m(k)}\)。于是，由备注 2 和例 2，有

\[
\mathrm{Q} _ {\mathrm{M}} = \operatorname{long} (\mathrm{H} _ {0}). \sum_ {0 \leqslant k \leqslant n} \sum_ {1 \leqslant j \leqslant m (k)} (- 1) ^ {k} \mathrm{T} ^ {\delta_ {k, j}},
\]

\[
c _ {\mathbf {M}} = \operatorname{long} (\mathrm{H} _ {0}). \sum_ {0 \leqslant k \leqslant n} (- 1) ^ {k} m (k).
\]

备注 4. --- 可以证明（p. 88, exerc. 4）：在定理 1 的假设下，\(\mathrm{H}_{0}\)-模 \(\mathrm{Tor}_{j}^{\mathrm{H}}(\mathrm{H}_{0},\mathbf{M})\) 具有有限长度，当 \(j > r\) 时为零，并且有 \(c_{\mathbf{M}} = \sum_{j=0}^{r} (-1)^{j} \mathrm{long}_{\mathrm{H}_{0}}(\mathrm{Tor}_{j}^{\mathrm{H}}(\mathrm{H}_{0},\mathbf{M}))\)。更确切地说，H-模

\(T_{j} = \text{Tor}_{j}^{\text{H}}(\text{H}_{0}, \text{M})\) 自然带有分次结构，并且有

\[
\mathrm{Q} _ {\mathrm{M}} = \sum_ {j = 0} ^ {r} (- 1) ^ {j} \mathrm{P} _ {\mathrm{T} _ {j}}.
\]

命题 1.--- 设 M 是分次 H-模。假设 M 由 \(M_{0}\) 生成，且 \(M_{0}\) 是具有有限长度的 \(H_{0}\)-模。则有

\[
\mathrm{P} _ {\mathrm{M}} \leqslant (1 - \mathrm{T}) ^ {- r} \operatorname{long} _ {\mathrm{H} _ {0}} (\mathrm{M} _ {0}), \quad c _ {\mathrm{M}} \leqslant \operatorname{long} _ {\mathrm{H} _ {0}} (\mathrm{M} _ {0}).
\]

此外，下列条件等价：

\begin{enumerate}
\def\labelenumi{(\roman{enumi})}
\tightlist
\item
  \(c_{\mathbf{M}} = \mathrm{long}_{\mathrm{H_0}}(\mathbf{M}_0)\)；
\end{enumerate}

\[
\text{(ii)} \mathrm{P}_{\mathbf{M}} = \operatorname{long}_{\mathbf{H}_{0}}(\mathbf{M}_{0}) \cdot (1 - \mathrm{T})^{-r},\ \text{即， }\mathbf{M} = \mathbf{M}_{0}\text{ 若 }r=0\text{ 且}
\]

\[
\operatorname{long} _ {\mathrm{H} _ {0}} \left(\mathrm{M} _ {n}\right) = \operatorname{long} _ {\mathrm{H} _ {0}} \left(\mathrm{M} _ {0}\right). \binom {n + r - 1} {r - 1}
\]

对 \(n \in \mathbf{N}\) 成立，当 \(r > 0\) 时；

\begin{enumerate}
\def\labelenumi{(\roman{enumi})}
\setcounter{enumi}{2}
\tightlist
\item
  H-模的典范同态
\end{enumerate}

\[
\varphi : \mathrm{H} \otimes_ {\mathrm{H} _ {0}} \mathrm{M} _ {0} \rightarrow \mathrm{M}
\]

是双射。

令 R 表示 φ 的核。由于 φ 是满射，有

\(P_{M} = P_{H \otimes M_{0}} - P_{R} = \text{long}_{H_{0}}(M_{0})(1 - T)^{-r} - P_{R}\) 且 \(c_{M} = \text{long}_{H_{0}}(M_{0}) - c_{R}\)。条件 (i)、(ii) 和 (iii) 分别等价于 \(c_{R} = 0\)、\(P_{R} = 0\) 和 \(R = 0\)。因此有 (iii) \(\Rightarrow\) (ii) \(\Rightarrow\) (i)，只需证明 \(c_{R} = 0\) 蕴含 \(R = 0\)。假设 \(R \neq 0\)，并令 \(0 = N^{h} \subset N^{h-1} \subset \ldots \subset N^{0} = M_{0}\) 为 \(H_{0}\)-模 \(M_{0}\) 的 Jordan–Hölder 序列。令 \(R^{m}\) 为 \(R\) 与 \(H \otimes_{H_{0}} N^{m}\) 在 \(H \otimes_{H_{0}} M_{0}\) 中的像的交；存在整数 \(m\)，它介于 1 和 \(h\) 之间，使 \(R^{m} \neq R^{m-1}\)。置 \(L = R^{m-1}/R^{m}\)；有 \(0 \leqslant c_{L} \leqslant c_{R}\)，只需证明 \(c_{L} \neq 0\)。现在，若 \(k\) 是 \(H_{0}\) 除以 \(N^{m-1}/N^{m}\) 的极大湮灭理想所得的商域，则 L 可识别为 \(k[(X_{i})_{i\in I}]\) 的非零分次子模。因此 L 包含一个与 \(k[(X_{i})_{i\in I}]\) 的平移模同构的子模；由于 \(c_{k[(X_{i})_{i\in I}]} = 1\)，所以 \(c_{L} \geqslant 1\)（备注 2 和 3），这正是我们要证明的。

根据 A, X, p. 160 的 th. 1，条件 (iii) 意味着 \((\mathbf{X}_{1}, \ldots, \mathbf{X}_{r})\) 是 H-模 M 的完全截切序列。

命题 2. --- 假设 \(H_{0}\) 是域，令 M 为有限生成的分次 H-模。令 K 为 H 的分式域。则 \(c_{M}\) 等于 H-模 M 的秩，即 K-向量空间 \(M \otimes_{H} K\) 的维数。

当 M = H 时这是显然的，因为 \(c_{H} = 1\)。此外，令 \(x \in H\) 为 d 次齐次且非零的元素；有 (H/xH) \(\otimes_{H} K = 0\)；由正合序列

\[
0 \rightarrow \mathrm{H} (- d) \rightarrow \mathrm{H} \rightarrow \mathrm{H} / x \mathrm{H} \rightarrow 0,
\]

以及备注 2 和 3，得到 \(c_{\mathrm{H}/x\mathrm{H}}=0\)。因此，当 M 由单个齐次元素生成时，命题成立。一般情形由此得到，因为每个有限生成分次 H-模都有一个合成序列，其商都是上述形式。

备注 6. -- 在 Prop. 2 的假设下，\(c_{\mathbf{M}} = 0\) 当且仅当 M 是挠 H-模，或者等价地，当且仅当 \(\dim_{\mathrm{H}}(\mathbf{M}) < r\)（\(\S 1, \mathfrak{n}^{\circ} 5\)，Example 4）。

